\documentclass[12pt,a4paper]{article}

% Encoding and fonts
\usepackage[utf8]{inputenc}
\usepackage[T1]{fontenc}
\usepackage{lmodern}
\usepackage{textcomp}

% Page layout
\usepackage[top=1in, bottom=1in, left=1in, right=1in]{geometry}

% Typography
\usepackage{microtype}
\usepackage{parskip}

% Language
\usepackage[english]{babel}

% Headers and footers
\usepackage{fancyhdr}
\pagestyle{fancy}
\fancyhf{}
\fancyhead[L]{\leftmark}
\fancyhead[R]{\thepage}
\fancyfoot[C]{\thepage}
\renewcommand{\headrulewidth}{0.4pt}
\renewcommand{\footrulewidth}{0pt}

% Section styling
\usepackage{titlesec}
\titleformat{\section}{\Large\bfseries}{\thesection}{1em}{}
\titleformat{\subsection}{\large\bfseries}{\thesubsection}{1em}{}
\titleformat{\subsubsection}{\normalsize\bfseries}{\thesubsubsection}{1em}{}
\titlespacing*{\section}{0pt}{2.5ex plus 1ex minus .2ex}{1.5ex plus .2ex}
\titlespacing*{\subsection}{0pt}{2ex plus 1ex minus .2ex}{1ex plus .2ex}
\titlespacing*{\subsubsection}{0pt}{1.5ex plus 1ex minus .2ex}{0.8ex plus .2ex}

% List formatting
\usepackage{enumitem}
\setlist[itemize]{topsep=0.5ex, itemsep=0.25ex, parsep=0pt}
\setlist[enumerate]{topsep=0.5ex, itemsep=0.25ex, parsep=0pt}

% Essential packages
\usepackage{graphicx}
\usepackage[table]{xcolor}
\usepackage{booktabs}
\usepackage{array}
\usepackage{tabularx}
\usepackage{longtable}
\usepackage{amsmath}
\usepackage{amssymb}
\usepackage[normalem]{ulem}
\rowcolors{2}{gray!10}{white}
\renewcommand{\arraystretch}{1.3}

% Cross-references
\usepackage{hyperref}
\usepackage{cleveref}

% Custom commands
\newcommand{\pilcrow}{\S}

% Hyperref setup
\hypersetup{
    colorlinks=true,
    linkcolor=blue,
    filecolor=magenta,
    urlcolor=cyan,
}

% Code listings
\usepackage{listings}
\definecolor{codebg}{RGB}{248,248,248}
\definecolor{codeframe}{RGB}{200,200,200}
\definecolor{codekeyword}{RGB}{0,0,180}
\definecolor{codecomment}{RGB}{0,128,0}
\definecolor{codestring}{RGB}{163,21,21}
\lstset{
    basicstyle=\ttfamily\small,
    breaklines=true,
    breakatwhitespace=false,
    frame=single,
    framerule=0.5pt,
    rulecolor=\color{codeframe},
    backgroundcolor=\color{codebg},
    keepspaces=true,
    numbers=left,
    numberstyle=\tiny\color{gray},
    numbersep=8pt,
    showstringspaces=false,
    tabsize=4,
    xleftmargin=1em,
    xrightmargin=1em,
    aboveskip=1em,
    belowskip=1em,
    keywordstyle=\color{codekeyword}\bfseries,
    commentstyle=\color{codecomment}\itshape,
    stringstyle=\color{codestring},
    literate={_}{{\textunderscore}}1
             {-}{{\textminus}}1
             {<}{{\textless}}1
             {>}{{\textgreater}}1
             {~}{{\textasciitilde}}1
             {^}{{\textasciicircum}}1
}

% YAML language definition for listings
\lstdefinelanguage{YAML}{
    keywords={true,false,null,y,n},
    keywordstyle=\color{blue}\bfseries,
    sensitive=false,
    comment=[l]{\#},
    morecomment=[s]{/*}{*/},
    commentstyle=\color{gray}\ttfamily,
    stringstyle=\color{red}\ttfamily,
    morestring=[b]'',
    morestring=[b]'',
}

% TOML language definition for listings
\lstdefinelanguage{TOML}{
    keywords={true,false},
    keywordstyle=\color{blue}\bfseries,
    sensitive=true,
    comment=[l]{\#},
    commentstyle=\color{gray}\ttfamily,
    stringstyle=\color{red}\ttfamily,
    morestring=[b]'',
    morestring=[b]'',
}

% Dockerfile language definition for listings
\lstdefinelanguage{Dockerfile}{
    keywords={FROM,RUN,CMD,LABEL,EXPOSE,ENV,ADD,COPY,ENTRYPOINT,VOLUME,USER,WORKDIR,ARG,ONBUILD,STOPSIGNAL,HEALTHCHECK,SHELL},
    keywordstyle=\color{blue}\bfseries,
    sensitive=false,
    comment=[l]{\#},
    commentstyle=\color{gray}\ttfamily,
    stringstyle=\color{red}\ttfamily,
    morestring=[b]'',
    morestring=[b]'',
}

% JSON language definition for listings
\lstdefinelanguage{JSON}{
    keywords={true,false,null},
    keywordstyle=\color{blue}\bfseries,
    sensitive=true,
    commentstyle=\color{gray}\ttfamily,
    stringstyle=\color{red}\ttfamily,
    morestring=[b]'',
}

% Code listing float environment
\usepackage{float}
\newfloat{listing}{htbp}{lol}[section]
\floatname{listing}{Listing}

% Admonition boxes
\usepackage{tcolorbox}
\tcbuselibrary{skins,breakable}

% Admonition style definitions
\newtcolorbox{admonitionbox}[2][]{
    enhanced,
    breakable,
    colback=#2!5,
    colframe=#2!80!black,
    fonttitle=\bfseries,
    title=#1,
    left=2mm,
    right=2mm,
}

% Synthon tuple boxes
\newtcolorbox{synthonbox}{
    enhanced,
    colback=blue!5,
    colframe=blue!75!black,
    fontupper=\large,
    center upper,
    boxrule=1pt,
    arc=4pt,
}

\begin{document}

\title{Frobenius-Guided Gene Editing: Perfecting Editing Techniques Through the Imscribing Grammar of the Genetic Code}
\date{\today}

\maketitle

\[\frac{17,280,000}{64} = 270,000\]

The codon space divides the Crystal of Types exactly. The genetic code is not a biological contingency --- it is a stratified Frobenius algebra on the B₄³ codon lattice, with $\mu \circ \delta = \text{id}$ holding exactly on the ground layer and up to $\mathbb{Z}_2$ wobble symmetry on the promoted layer. This structural fact constrains what gene editing \textbf{is} --- and therefore how it can be made \textbf{perfect}.

\begin{center}
\rule{0.5\textwidth}{0.4pt}
\end{center}

\subsection{The Editing Problem in Structural Terms}

Every gene editing operation --- CRISPR-Cas9 double-strand break, base editing, prime editing, homology-directed repair --- performs a \textbf{local modification of the Frobenius algebra on codon space}. The edit target is a locus in B₄³. The edit outcome is a new codon (or set of codons) with a new position in the Frobenius stratification: exact stratum, split stratum, or stop stratum. The structural cost of an edit is the Frobenius distance between the original and target codons.

Current editing techniques share a common limitation: they optimize for \textbf{sequence homology} (Watson-Crick base pairing) but not for \textbf{Frobenius closure}. A sequence-perfect edit can still be structurally defective if it disrupts the stratified Frobenius structure of the genetic code in ways that propagate through translation, folding, and function.

The standard approach asks: ''''Does this edit change the amino acid?'''' The IG approach asks: ''''Does this edit change the Frobenius stratum of the codon, and if so, what primitive activation does it alter at the protein''s structural type?''''

\begin{center}
\rule{0.5\textwidth}{0.4pt}
\end{center}

\subsection{B₄ Lattice Topology of Nucleotide Editing}

The B₄ lattice (G $\rightarrow$ B, C $\rightarrow$ T, A $\rightarrow$ F, U $\rightarrow$ N) governs the structural relationships between nucleotides:

\begin{lstlisting}
        B (Both)
       / \
      T   N
       \ /
        F
\end{lstlisting}

\textbf{Covering relations} (minimal structural transitions):

\begin{itemize}
    \item B $\rightarrow$ T, B $\rightarrow$ N (Both $\rightarrow$ True, Both $\rightarrow$ Neither)
    \item T $\rightarrow$ F, N $\rightarrow$ F (True $\rightarrow$ False, Neither $\rightarrow$ False)
\end{itemize}

\textbf{Non-covering transitions} (cross-lattice jumps):

\begin{itemize}
    \item T $\leftrightarrow$ N (True $\leftrightarrow$ Neither)
    \item B $\leftrightarrow$ F (Both $\leftrightarrow$ False)
    \item G$\leftrightarrow$A, G$\leftrightarrow$U, C$\leftrightarrow$A, C$\leftrightarrow$U (direct purine/pyrimidine swaps)
\end{itemize}

\textbf{Structural cost of base edits:}

\begin{table}[htbp]
\centering
\small
\rowcolors{1}{white}{white}
\begin{tabularx}{\textwidth}{>{\centering\arraybackslash}X >{\centering\arraybackslash}X >{\centering\arraybackslash}X >{\centering\arraybackslash}X}
\toprule
\rowcolor{gray!25}\textbf{Edit} & \textbf{B₄ transition} & \textbf{Lattice type} & \textbf{Structural cost} \\
\midrule
\rowcolors{1}{white}{gray!10}
C$\rightarrow$T (transition) & T $\rightarrow$ F & Covering & Minimal \\
G$\rightarrow$A (transition) & B $\rightarrow$ F & Non-covering (cross-lattice) & Maximal \\
C$\rightarrow$A (transversion) & T $\rightarrow$ ? & Crossing & Intermediate \\
G$\rightarrow$U (transversion) & B $\rightarrow$ N & Covering & Minimal \\
A$\rightarrow$G (transition) & F $\rightarrow$ B & Non-covering (cross-lattice) & Maximal \\
U$\rightarrow$C (transition) & N $\rightarrow$ T & Covering & Minimal \\
\bottomrule
\end{tabularx}
\end{table}

\textbf{Base editors} (CBE, ABE) are limited to transitions. The B₄ lattice reveals why this is not merely a biochemical limitation of deaminase enzymes --- it is a \textbf{structural constraint}: covering relations preserve the B₄ hierarchy. Transitions on the covering path (T$\leftrightarrow$F, N$\leftrightarrow$T, F$\leftrightarrow$B, B$\leftrightarrow$N) are lattice-coherent. Cross-lattice transitions (B$\leftrightarrow$F, T$\leftrightarrow$N) are structurally discontinuous.

\textbf{Key insight:} A perfect base editor is one that operates only along B₄ covering relations. ABE (A$\rightarrow$G, F$\rightarrow$B) necessarily performs a \textbf{maximal structural jump} --- from False to Both --- which is why ABE editing of A$\cdot$T pairs at the wobble position of a split-stratum codon can produce unanticipated amino acid outcomes. The error is not thermodynamic; it is lattice-topological.

\begin{center}
\rule{0.5\textwidth}{0.4pt}
\end{center}

\textbf{Author:} Lando \otimes ⊙perator

\subsection{Frobenius Stratum and Guide RNA Design}

The 16 codon boxes split 8/8 into Frobenius-exact and Frobenius-open strata. The rule: \textbf{exact iff p₂ = T (C at position 2), or p₂ $\in$ \{N, B\} (G/U) with p₁ $\in$ \{T, B\} (C/G).}

\subsubsection{Guide RNA targeting within the exact stratum}

When a guide RNA targets a codon in one of the 8 exact boxes (32 codons), \textbf{position 3 carries no information}. Any edit at position 3 is silent at the amino acid level. However, a guide that cleaves at position 3 of an exact-stratum codon necessarily triggers repair across a Frobenius-closed region --- the repair template must also preserve the exactness condition, otherwise the edited codon will belong to a different box entirely.

\textbf{Design rule:} For exact-stratum targets, the guide RNA should span positions 1--2 of the codon, not position 3. The PAM-proximal seed region should be positioned to distinguish between adjacent exact and split boxes --- a single-nucleotide shift can move the target from the exact box CU\_ (Leu, position 2 = C = T) to the split box UU\_ (Phe/Leu, position 2 = U = N). The Frobenius distance of such a shift is the distance between T and N in B₄ --- a non-covering transition --- which is structurally maximal.

\subsubsection{Guide RNA targeting within the split stratum}

For split-stratum targets (8 boxes, 29 codons), position 3 carries the pyrimidine/purine distinction. A guide that cleaves at position 3 of a split-stratum codon \textbf{must} distinguish pyrimidine (Y) from purine (R), but need not distinguish the four individual bases.

\textbf{Design rule:} Split-stratum guides should terminate at position 3 with wobble-tolerant bases (inosine, G$\cdot$U wobble pairs) that respect the pyrimidine/purine equivalence relation. Current gRNA design treats all four bases as independent --- this over-constrains the search space by a factor of 2 at every split-stratum position. The Frobenius structure says a 2-fold reduction is legitimate.

\subsubsection{Crossing the exact$\leftrightarrow$split boundary}

An edit that moves a codon from an exact box to a split box (or vice versa) is a \textbf{Frobenius stratum crossing}. This is the editing equivalent of a topological phase transition. The consequences:

\begin{itemize}
    \item \textbf{Exact $\rightarrow$ Split:} Position 3 gains informational content. A previously silent site becomes meaningful. If the cell''s repair machinery fills position 3 with a nucleotide that changes the pyrimidine/purine state, the amino acid changes. This is the structural source of many Cas9 off-target effects --- the guide targets the exact stratum, but the repair machinery interprets position 3 according to split-stratum rules.
    \item \textbf{Split $\rightarrow$ Exact:} Position 3 loses informational content. Two previously distinct amino acids (e.g., His/His from CA\_ pyrimidine/purine) collapse to one. The repair template need not specify position 3 at all --- a degenerate base (N) at that position is Frobenius-optimal.
\end{itemize}

\textbf{Practical consequence:} A perfect editing protocol should never design a guide whose on-target site is in the exact stratum but whose off-target sites include split-stratum codons, or vice versa. Current off-target prediction algorithms score by sequence similarity alone --- they miss this categorical difference.

\begin{center}
\rule{0.5\textwidth}{0.4pt}
\end{center}

\subsection{The 12 Primitive Activations as Editing Constraints}

Each of the 12 promoted amino acids instantiates exactly one IG primitive that the ground layer (8 exact-box amino acids) does not activate. When editing a coding sequence, changing one promoted amino acid to another changes which primitive is active at that protein position.

\subsubsection{Primitive load-bearing map}

\begin{table}[htbp]
\centering
\small
\rowcolors{1}{white}{white}
\begin{tabularx}{\textwidth}{>{\centering\arraybackslash}X >{\centering\arraybackslash}X >{\centering\arraybackslash}X}
\toprule
\rowcolor{gray!25}\textbf{AA} & \textbf{Primitive} & \textbf{Editing constraint} \\
\midrule
\rowcolors{1}{white}{gray!10}
Met & Ð (scope) & Editing the start codon (AUG) changes the entire translation scope. The only edit that preserves Frobenius closure is AUG $\rightarrow$ AUG (silent at position 3) or AUG $\rightarrow$ alternative start (rare, constrained). \\
Trp & Þ (topology) & The bicyclic indole is the topological ceiling. Any Trp $\rightarrow$ X edit collapses the topological complexity of that residue to a lower tier. \\
Cys & Ř (reversibility) & Disulfide bonds are the only reversible covalent crosslinks in proteins. Editing Cys out of a structurally critical disulfide pair breaks the reversible relation --- permanent structural drift. \\
Tyr & $\Phi$ (parity) & The +OH group enables phosphorylation (parity switch). Editing Tyr $\rightarrow$ Phe removes the switch; editing Tyr $\rightarrow$ other aromatics may preserve volume but not parity. \\
Phe & ƒ (force ceiling) & Maximum hydrophobicity. Editing Phe $\rightarrow$ Leu or Ile preserves hydrophobicity partially; editing to a polar residue is a fidelity collapse. \\
Ile & Ç (kinetic constraint) & $\beta$-branched stereocenter generates the tightest ribosomal decoding coupling. Editing Ile $\rightarrow$ Val preserves $\beta$-branching; editing to anything else relaxes the kinetic constraint. \\
His & $\Gamma$ (grammar switch) & Imidazole pKa $\approx$ 6 bridges acid/base catalysis. Editing His removes the pH-gated catalytic grammar --- unless replaced by another amphoteric side chain (none exist among the 20). \\
Asn & ɢ (interaction) & N-glycosylation site (N-X-S/T sequon). Editing Asn removes the extracellular recognition gate. \\
Gln & ⊙ (criticality) & Glutamine synthetase is the most regulated biosynthetic node in metabolism. Editing Gln at regulatory positions can trigger runaway (supercritical) metabolism. \\
Asp & Ħ (chirality) & Asp in active sites enforces chiral substrate selectivity. Editing Asp $\rightarrow$ Glu preserves length change but not chiral specificity. \\
Lys & $\Sigma$ (entropy) & Highest sequence variability + epigenetic acetylation target. Editing Lys $\rightarrow$ Arg preserves charge but loses the acetylation switch (methylation differential). \\
Glu & $\Omega$ (winding) & Highest $\alpha$-helix propensity; helix dipole stabilizer. Editing Glu at helix N-termini disrupts helix winding initiation. \\
\bottomrule
\end{tabularx}
\end{table}

\subsubsection{Editing risk by primitive class}

\textbf{Frobenius-locked primitives (editing at these positions is catastrophic):}

\begin{itemize}
    \item \textbf{Ħ (chirality):} All 19 chiral amino acids are L-configuration. D-amino acid insertion at any position breaks the ribosomal Frobenius gate. Base editing cannot produce D-amino acids (the chirality is post-translational), but any edit that inserts a non-standard amino acid capable of racemization at the active site is structurally risky.
    \item \textbf{Ð (scope):} The start codon defines translation scope. No legitimate edit can delete Met without destroying the reading frame or initiating at a non-canonical start.
    \item \textbf{$\Omega$ (winding):} Stop codons are the $\Omega$ closure signal. Editing a stop codon to a sense codon (readthrough) is equivalent to removing the winding boundary. This is permitted in rare cases (e.g., selenocysteine recoding of UGA) but requires the full $\Omega$\_Z topological protection mechanism.
\end{itemize}

\textbf{Frobenius-flexible primitives (editing at these positions is least disruptive):}

\begin{itemize}
    \item \textbf{$\Sigma$ (entropy):} Lys is the most variable amino acid by sequence conservation metrics. Lys $\rightarrow$ Arg is the most common conserved substitution across the proteome --- it preserves charge but changes acetylation potential. This is the lowest-risk mutation class in the promoted layer.
    \item \textbf{ƒ (force):} Hydrophobic substitutions within the same packing class (Phe $\leftrightarrow$ Leu $\leftrightarrow$ Ile $\leftrightarrow$ Val) preserve the fidelity regime. The risk gradient is monotonic with hydrophobicity difference.
\end{itemize}

\textbf{Semi-locked primitives (editing requires compensatory changes):}

\begin{itemize}
    \item \textbf{Ř (reversibility):} Editing a single Cys out of a disulfide pair breaks the reversible linkage. If the edit is necessary (e.g., for therapeutic protein stability engineering), the partner Cys must also be edited or the disulfide bond becomes an orphan half-cystine --- a Frobenius-open defect that can cause aggregation.
    \item \textbf{⊙ (criticality):} Editing Gln at a regulatory position changes the criticality of the metabolic node. A compensatory edit elsewhere in the same pathway may restore the critical point.
\end{itemize}

\begin{center}
\rule{0.5\textwidth}{0.4pt}
\end{center}

\subsection{Prime Editing and the Frobenius Template Rule}

Prime editing is the closest existing technique to a Frobenius-optimal editing method. The prime editing guide RNA (pegRNA) encodes both the nick site and the edit template --- a $\delta$ (comultiplication) that specifies the target codon. The repair product is the $\mu$ (multiplication) of that template. Prime editing succeeds when $\mu$∘$\delta$ = id for the edited locus.

\subsubsection{The Frobenius template rule}

For any prime edit, the structural condition for a ''''perfect'''' edit is:

\[\mu(\delta(\text{target codon})) = \text{target codon}''\]

where $\mu$ is the genetic code table, $\delta$ is the pegRNA-specified template, and the composite $\mu$∘$\delta$ must equal identity on the \textbf{Frobenius stratum} of the target. This decomposes into three conditions:

\begin{enumerate}
    \item \textbf{Stratum preservation:} The edit must not move the target codon from exact to split or split to exact unless explicitly intended.
    \item \textbf{Primitive invariance:} If the edit changes the amino acid, the new amino acid should activate the same IG primitive as the original, or the protein should be evolutionarily tolerant of that primitive swap.
    \item \textbf{$\Omega$ boundary respect:} If the edit is near a stop codon, the stop''s $\Omega$ value must be preserved or the protein''s C-terminal winding must be redesigned.
\end{enumerate}

\subsubsection{pegRNA design improvements from the IG grammar}

Current pegRNA design optimizes for:

\begin{itemize}
    \item Reverse transcriptase template length (10--15 nt)
    \item Primer binding site length (8--13 nt)
    \item GC content
    \item Homology arm symmetry
\end{itemize}

The IG grammar adds three new optimization criteria:

\begin{itemize}
    \item \textbf{Frobenius stratum of edit window:} Is the target site in an exact or split box? Position 3 edits in exact boxes are silent (low risk); position 3 edits in split boxes change amino acids (high risk).
    \item \textbf{B₄ lattice distance of edit:} What is the covering relation between the original and target nucleotide? Edits along B₄ covering relations (lattice distance = 1) are structurally minimal.
    \item \textbf{Primitive load of product amino acid:} Does the edit create an amino acid from a different primitive class? If so, what is the structural distance between the primitives?
\end{itemize}

\textbf{Example --- correcting a pathogenic G$\rightarrow$A mutation:}
A G$\rightarrow$A mutation at codon position 1 is a B $\rightarrow$ F transition (maximal structural cost, cross-lattice). If the target is in the split stratum, the pegRNA must specify all three positions of the edited codon. Current pegRNA design would use a 10-nt RT template centered on the edit site. The IG-aware design would extend the RT template to include the neighboring codon''s position 2 (which determines the Frobenius stratum) and add degenerate bases at position 3 if the target is in the exact stratum.

\begin{center}
\rule{0.5\textwidth}{0.4pt}
\end{center}

\subsection{Homology-Directed Repair and the Sheaf Structure}

HDR templates are the $\delta$ map of the Frobenius algebra --- they specify the intended edit. The cellular repair machinery executes $\mu$ by integrating the template. The sheaf structure of the genetic code (the property that meet-closure holds within each Frobenius stratum but not across strata) constrains HDR template design.

\subsubsection{Sheaf-optimal HDR templates}

An HDR template is sheaf-optimal when it:

\begin{enumerate}
    \item \textbf{Preserves the local Frobenius stratum.} If the target locus spans multiple codon boxes, the template must specify all positions to maintain stratum membership.
    \item \textbf{Does not create orphan Frobenius defects.} An edit that changes one amino acid but leaves the neighboring codons in a different stratum creates a boundary defect --- the ribosomal $\mu$ map sees a discontinuity.
    \item \textbf{Uses degenerate bases at Frobenius-closed positions.} In exact-stratum codons, position 3 can be specified as N (any base). This increases the probability of successful HDR by 4$\times$ without increasing off-target amino acid risk.
\end{enumerate}

\subsubsection{The Cas9 off-target sheaf theorem}

\textbf{Theorem:} For any CRISPR guide RNA with off-target sites in a different Frobenius stratum than the on-target site, the probability of a structural defect at the off-target site is at least 50\% (the probability that repair fills position 3 with the ''''wrong'''' pyrimidine/purine).

\textbf{Proof sketch:} On-target site is in stratum S. Off-target site is in stratum S'' $\neq$ S. The cellular repair machinery responds to the double-strand break with NHEJ or HDR. If HDR uses the on-target template, the off-target site''s Frobenius stratum is mismatched --- position 3 is filled according to the on-target stratum rules, which are incorrect for S''. The probability that the random filling of position 3 happens to match the correct pyrimidine/purine for S'' is 1/2. ■

This theorem explains a known but poorly understood empirical pattern: Cas9 off-target edits are enriched for amino acid changes at position 3 of split-stratum codons. The structural mechanism is not sequence homology; it is Frobenius stratum mismatch.

\subsection{Therapeutic Editing and Structural Risk Stratification}

The 12-primitive activation map enables a risk stratification system for therapeutic gene editing that is independent of any particular delivery or editing modality.

\subsubsection{Positional risk by primitive class}

For any coding-sequence edit, compute the \textbf{primitive delta} between the original and edited amino acids:

\begin{table}[htbp]
\centering
\small
\rowcolors{1}{white}{white}
\begin{tabularx}{\textwidth}{>{\centering\arraybackslash}X >{\centering\arraybackslash}X >{\centering\arraybackslash}X}
\toprule
\rowcolor{gray!25}\textbf{Primitive change} & \textbf{Risk class} & \textbf{Example therapeutic scenario} \\
\midrule
\rowcolors{1}{white}{gray!10}
Ħ $\rightarrow$ any & \textbf{Critical} & Editing an Asp (chirality-enforcing) at an active site. Any replacement loses chiral specificity. \\
Ð $\rightarrow$ any & \textbf{Critical} & Editing the start codon to anything other than AUG. Translational scope destroyed. \\
$\Omega$ $\rightarrow$ any & \textbf{Critical} & Editing a stop codon to a sense codon without the selenocysteine machinery. Readthrough produces C-terminal extension. \\
Ř $\rightarrow$ non-Ř & \textbf{High} & Editing a disulfide-forming Cys. Requires partner Cys edit --- compensatory pair design needed. \\
⊙ $\rightarrow$ non-⊙ & \textbf{High} & Editing Gln at a regulatory metabolic node. Requires metabolic criticality analysis. \\
Þ $\rightarrow$ non-Þ & \textbf{Moderate} & Trp $\rightarrow$ any. Indole collapse tolerated in surface positions but not in core packing. \\
$\Phi$ $\rightarrow$ non-$\Phi$ & \textbf{Moderate} & Tyr $\rightarrow$ Phe (loss of phosphorylation site). Functional impact depends on whether the site is a known signaling switch. \\
Ç $\rightarrow$ non-Ç & \textbf{Moderate} & Ile $\rightarrow$ Val ($\beta$-branching preserved) vs Ile $\rightarrow$ Leu ($\beta$-branching lost). The latter is higher risk. \\
$\Gamma$ $\rightarrow$ non-$\Gamma$ & \textbf{Moderate} & His $\rightarrow$ any at active site pH gate. Requires redesign of catalytic mechanism. \\
ɢ $\rightarrow$ non-ɢ & \textbf{Moderate} & Asn $\rightarrow$ any at N-glycosylation sequon. Loss of glycosylation is often pathological. \\
$\Sigma$ $\rightarrow$ $\Sigma$ & \textbf{Low} & Lys $\rightarrow$ Arg (charge preserved, acetylation lost vs methylation gained). Most tolerated. \\
ƒ $\rightarrow$ ƒ (within class) & \textbf{Low} & Phe $\leftrightarrow$ Leu $\leftrightarrow$ Ile $\leftrightarrow$ Val. Hydrophobic core packing adjustments. \\
Any ground $\rightarrow$ ground & \textbf{Low} & Gly, Ala, Pro, Ser, Thr, Val (exact-box AAs) $\rightarrow$ any other ground AA. These have no primitive activation --- any ground substitution preserves the Frobenius-exact stratum structure. \\
\bottomrule
\end{tabularx}
\end{table}

\subsubsection{The Chimera Theorem}

When editing across primitive classes, the risk is not additive --- it is \textbf{tensorial}. The composite risk of a double edit that changes two primitives is the tensor product of the individual primitive risks, not their sum.

\textbf{Example:} Editing Cys (Ř) at a disulfide bridge AND His ($\Gamma$) at the same active site creates a composite risk = \otimes(Ř, $\Gamma$). The tensor product of two semi-locked primitives is a \textbf{trap state} (Ç\_⊛): once both edits are made, the protein may be structurally frozen in a non-functional conformation from which no further editing can rescue it. Single edits on either primitive alone are recoverable.

\textbf{Clinical implication:} Therapeutic editing protocols should be screened for tensor products that produce trap states. A single observed off-target edit may be tolerable; two off-target edits at distinct primitive classes may be structurally catastrophic regardless of their genomic distance.

\begin{center}
\rule{0.5\textwidth}{0.4pt}
\end{center}

\subsection{Beyond CRISPR: IG-Optimal Editing Architecture}

The IG grammar suggests a novel editing architecture that is not limited by current biochemical implementations:

\subsubsection{Frobenius-Respecting Nuclease}

Current nucleases (Cas9, Cas12, Cas13) cleave based on RNA-guided sequence recognition. A Frobenius-optimal nuclease would:

\begin{enumerate}
    \item \textbf{Stratum-classify the target} before cleavage: is the target in the exact or split stratum?
    \item \textbf{Select cleavage position by stratum:} For exact-stratum targets, cleave at positions 1--2 only (leave position 3 as N). For split-stratum targets, cleave at all three positions with pyrimidine/purine degeneracy at position 3.
    \item \textbf{Verify Frobenius closure post-cleavage:} Confirm that the repair product satisfies $\mu$∘$\delta$=id for the target locus.
\end{enumerate}

\subsubsection{Multi-primitive editing}

Current editing is single-locus. The sheaf structure of the genetic code means that editing a single codon can have structural effects on neighboring codons --- particularly at stratum boundaries. An IG-optimal editing system would:

\begin{enumerate}
    \item \textbf{Edit in primitive-compatible groups} --- not single nucleotides but sets of positions that together preserve the Frobenius algebra.
    \item \textbf{Use the Crystal of Types} (17,280,000 addresses) to pre-compute the Frobenius-optimal edit path from any source codon to any target codon, including the minimal B₄ lattice distance and the risk of unintended stratum crossing.
    \item \textbf{Provide a Frobenius confidence score} for each edit, analogous to the C-score (consciousness score) but on the genetic code''s Frobenius algebra instead of the observer''s self-modeling loop.
\end{enumerate}

\subsubsection{The Editing Compiler}

The whale engine architecture (whale\_engine.py) provides a template: compile a desired protein edit (target amino acid change) through IMASM instructions into a Frobenius-optimal editing protocol. The three-stage pipeline is:

\begin{lstlisting}
Desired AA change
  -> Amino acid substitution -> Codon change -> Nucleotide edit
  -> Frobenius stratum check -> B₄ lattice path
  -> Guide RNA design -> Repair template design
  -> Frobenius closure verification -> Risk score
\end{lstlisting}

Each stage is a deterministic computation on the Crystal of Types. No empirical trial-and-error is needed for the structural part --- only for the biochemical implementation.

\begin{center}
\rule{0.5\textwidth}{0.4pt}
\end{center}

\subsection{Open Questions and Testable Predictions}

\subsubsection{Predictions}

\begin{enumerate}
    \item \textbf{Off-target structural bias:} Cas9 off-target edits in coding regions are enriched for exact$\leftrightarrow$split stratum crossings at a rate significantly above random (p \textless{} 0.01). This can be tested by re-analyzing published GUIDE-seq datasets with Frobenius stratum annotation.
    \item \textbf{Base editing fidelity by B₄ distance:} Adenine base editors (A$\rightarrow$G) have higher structural error rates at split-stratum position 3 than at exact-stratum position 3, because the A$\rightarrow$G (F$\rightarrow$B) transition at a split-stratum position 3 changes the pyrimidine/purine assignment. This is a testable prediction --- compare ABE editing outcomes at split-stratum vs exact-stratum targets with matched sequence context.
    \item \textbf{PE efficiency by primitive distance:} Prime editing efficiency correlates inversely with the primitive delta between original and edited amino acids. Edits within the same primitive class have higher efficiency than edits that cross primitive classes, independent of sequence context.
\end{enumerate}

\subsubsection{Questions}

\begin{enumerate}
    \item \textbf{Can the Frobenius condition be used to design Cas9 variants with stratum-specific PAM recognition?} A Cas9 variant that prefers NGG PAMs in exact strata and NGN PAMs in split strata would be a Frobenius-optimized tool.
    \item \textbf{Does the sheaf structure predict synthetic lethality?} If two coding edits are in different Frobenius strata that share a common meet, the meet is a synthetic lethal target --- editing either individually is safe, editing both is lethal.
    \item \textbf{Is the 12-primitive map of amino acids the only possible map?} If the genetic code had evolved differently (different codon assignments), would the same 12-primitive activation pattern re-emerge? The IG grammar predicts yes --- the 12 primitives are the complete set of structural dimensions, and any functional genetic code must instantiate all 12.
\end{enumerate}

\begin{center}
\rule{0.5\textwidth}{0.4pt}
\end{center}

\subsection{Conclusion}

Gene editing is currently practiced as a sequence-level technology. The Imscribing Grammar reveals that it is actually a \textbf{Frobenius algebra operation} on a stratified lattice with 12 primitive activation channels. The implications are not metaphorical --- they produce testable predictions about off-target effects, base editing fidelity, prime editing efficiency, and therapeutic risk stratification.

A perfected editing technique is one that:

\begin{itemize}
    \item Respects the B₄ lattice topology of nucleotide transitions
    \item Preserves the Frobenius stratum of the target locus
    \item Accounts for the primitive activation class of each amino acid
    \item Uses degenerate bases at Frobenius-closed positions
    \item Verifies $\mu$∘$\delta$=id for the edited product
\end{itemize}

The genetic code is not a frozen accident. It is a Frobenius algebra. Editing it should be treated as such.

\begin{center}
\rule{0.5\textwidth}{0.4pt}
\end{center}

\textbf{Author:} Lando \otimes ⊙perator


\end{document}
